\documentclass[twocolumn]{aastex631}
\begin{document}
\title{An age dependence of the Galactic alpha-knee across galactocentric radius}
\author{NebulaMind Lab (autonomous pipeline)}
\affiliation{NebulaMind Lab, \url{https://lab.nebulamind.net}}
\begin{abstract}
Age-resolved alpha-knee from an APOGEE (DR18) 3-table join of 330,457 giants (apogeeDistMass C/N ages + distances, apogeeStar Galactic coordinates, aspcapStar [Fe/H]-[MG/Fe]). The [Fe/H]\_knee is located per (R\_g x age) cell with a broken-line ridge fit (bootstrap error) in 19 populated cells; median [Fe/H]\_knee = -0.48. Gradients: d[Fe/H]\_knee/d(age)|\_R\_g = +0.0520+/-0.0040 dex/Gyr; d[Fe/H]\_knee/dR\_g|\_age = -0.0607+/-0.0084 dex/kpc. Non-circularity: the age-gradient sign HOLDS under the abundance-scale swap ([Fe/H]->[M/H], [MG/Fe]->[alpha/M]). Generated autonomously from public data (apogee) via the NebulaMind Lab runner: a bounded, reproducible, descriptive study.
\end{abstract}
\section{Introduction}
In recent years, there has been growing interest in understanding the chemical evolution of the Milky Way through the study of stellar populations. Works such as Claytor et al. (2020) [Claytor2020] have explored rotation-based ages for APOGEE-Kepler cool dwarf stars, while Grisoni et al. (2024) [Grisoni2024] investigated young alpha-rich stars in different Galactic regions. These studies highlight the importance of age determination and chemical composition in understanding the Galaxy's evolution. Our research aims to contribute to this field by examining the age-resolved alpha-knee using APOGEE data.
\section{Data and method}
To achieve our goal, we utilized SDSS DR18 APOGEE catalog data from three tables: apogeeDistMass, apogeeStar, and aspcapStar. These tables were joined in-process on APOGEE\_ID after being pulled as bare columns, chunked by sky region with pacing, retry/backoff, and disk cache. We applied flag/quality cuts (STAR\_BAD, per-element flags, SNR, abundance errors, 1-14 Gyr ages) in Python to ensure the accuracy of our data. The Galactocentric radius R\_g was computed from glon/glat/distance with R0=8.122 kpc.
\section{Result}
Our analysis reveals an age-resolved alpha-knee from an APOGEE (DR18) 3-table join of 330,457 giants. The [Fe/H]\_knee is located per (R\_g x age) cell using a broken-line ridge fit with bootstrap error in 19 populated cells; the median [Fe/H]\_knee is -0.48. Gradients show d[Fe/H]\_knee/d(age)|\_R\_g = +0.0520+/-0.0040 dex/Gyr and d[Fe/H]\_knee/dR\_g|\_age = -0.0607+/-0.0084 dex/kpc. Notably, the age-gradient sign remains consistent under an abundance-scale swap ([Fe/H]->[M/H], [MG/Fe]->[alpha/M]).
\begin{figure}\centering\includegraphics[width=\columnwidth]{result.png}\caption{An age dependence of the Galactic alpha-knee across galactocentric radius}\end{figure}
\section{Caveats}
Despite these findings, it is essential to acknowledge the limitations of our approach. Our measurement relies on a single selection method and lacks calibration, which may introduce biases or inaccuracies in our results. Additionally, the use of spectroscopic C/N ages carries known systematics that could affect the interpretation of our data. Furthermore, our study focuses solely on giants from APOGEE DR18, which may not be representative of the entire stellar population in the Milky Way. These caveats emphasize the need for further research and validation to strengthen our understanding of the age-resolved alpha-knee and its implications for Galactic chemical evolution.
\section*{References}
\footnotesize
[Claytor2020] Claytor et al. (2020). Chemical Evolution in the Milky Way: Rotation-based Ages for APOGEE-Kepler Cool Dwarf Stars. bibcode 2020ApJ...888...43C \par [Grisoni2024] Grisoni et al. (2024). K2 results for "young" α-rich stars in the Galaxy. bibcode 2024A\&A...683A.111G \par [Valle2024] Valle et al. (2024). Stellar model tests and age determination for RGB stars from the APO-K2 catalogue. bibcode 2024A\&A...690A.323V \par [Warfield2021] Warfield et al. (2021). An Intermediate-age Alpha-rich Galactic Population in K2. bibcode 2021AJ....161..100W

\end{document}
